
\subsubsection{3.1 Problem Formulation}

Bidirectional alignment is formulated as multi-objective RLHF:

$\max_\theta \mathbb{E}_{x \sim D, y \sim \pi_\theta(\cdot|x)} [R_{\text{combined}}(x, y) - \beta_{\text{KL}} \cdot D_{\text{KL}}(\pi_\theta \| \pi_{\text{ref}})]$

where:

$R_{\text{combined}}(x, y) = \alpha \cdot R_{\text{help}}(x, y) + \beta \cdot R_{\text{ctrl}}(x, y)$

\subsubsection{3.2 IFEval as Differentiable Training Signal}

IFEval defines 25 verifiable constraint types. Binary constraint checks are non-differentiable; they are converted to continuous scores via sigmoid soft thresholds:

$s_i = \sigma((v_i - t_i) / \tau)$

where $v_i$ is the observed value, $t_i$ is the threshold, and $\tau$ is a temperature parameter controlling smoothness.

The controllability reward aggregates soft constraint scores:

$R_{\text{ctrl}}(x, y) = \frac{1}{|C(x)|} \cdot \sum_{c_i \in C(x)} s_i$

Gradient flow was validated in experiment H-E1: variance = 0.039, scale.grad = 274.12, confirming non-trivial signal with gradient capability.

\subsubsection{3.3 Experimental Conditions}

\textbf{Treatments:} T1 ($\alpha =0.2, \beta =0.8)$, T2 ($\alpha =0.4, \beta =0.6)$, T3 ($\alpha =0.6, \beta =0.4)$, T4 ($\alpha =0.8, \beta =0.2)$

\textbf{Baselines:}
\begin{itemize}
\item B1: SFT-only (no RLHF)
\item B2: Helpfulness RLHF ($\alpha =1.0, \beta =0.0)$
\item B3: Quality RLHF (alternative reward)
\end{itemize}

\textbf{Data Split:} IFEval 70/30 (training/held-out), with held-out evaluation on 162 prompts.
